\DIFbluesubsection{RQ6: How strong are \toolname's specifications versus the baselines?}
\label{subsec:rq6-strength}
\DIFaddbegin
Validity certifies that a specification is discharged by its verifier, but it does not rank how informative the specification is. We therefore evaluate specification strength on the both-valid intersection using two complementary methods: an LLM-based semantic judge as the primary oracle, and a Frama-C/WP-based implication check as supporting evidence on the subset that SMT solvers can adjudicate.
\DIFaddend



\DIFaddbegin
\smallskip
\noindent\textbf{LLM-as-a-judge evaluation.}
Table~\ref{tab:strength_summary} reports pairwise strength judgments on the both-valid pairs. Pooled across models, \toolname\ is judged stronger than \autospec\ on 50.0\% of pairs, against 31.5\% in the reverse direction; against \specgen\ the comparison is closer, with \toolname\ ahead on 36.1\% and \specgen\ on 34.3\%. These numbers should be read together with each tool's mechanism rather than as a single victory metric, because the three tools obtain strength in fundamentally different ways.
\DIFaddend

\DIFaddbegin
\autospec\ instead compensates with exhaustive accumulation: it repeatedly samples the LLM to grow a large, redundant pool of candidate \texttt{ensures} and \texttt{loop invariant} clauses and filters by WP-validity, so the load-bearing facts are typically buried among many tautological or trivially implied ones; its strength comes from brute-force sampling rather than from semantic reasoning. \toolname\ takes a third route: at every generation boundary it injects the local symbolic context to steer the LLM toward more path-sensitive, non-redundant clauses than the \autospec\ sampler reliably reaches. On loop-free programs that stay within the flattened-memory model, the symbolic post-state is covered by Proposition~\ref{thm:sp-sound-complete}: within the same flattened-memory state-pair vocabulary, it entails any postcondition over the covered locations.
\DIFaddend

\DIFaddbegin
The model-wise results are not uniformly in \toolname's favor. On \texttt{gpt-5-mini}, \autospec\ wins more \autospec--\toolname\ pairs than \toolname\ does (45.1\% vs.\ 34.3\%). We attribute this to weaker models being less reliable at instantiating the symbolic relations exposed by \toolname: invalid or over-strong clauses are rejected by WP and then weakened or removed during refinement, reducing final specification strength. On stronger models the direction reverses, with \toolname\ stronger on 47.9\% of pairs for \texttt{gpt-5.4-mini} and 58.9\% for \texttt{gpt-5}, suggesting that symbolic guidance scales with model capability.
\DIFaddend

\DIFaddbegin
For strength comparison on the both-valid pairs, we employ an LLM-based pairwise judge (Fig.~\ref{fig:prompt-judge}), which evaluates semantic strength rather than verifier-specific acceptance and explicitly ignores language-specific syntax. It nonetheless inherits the limitations of LLM-based judgment, which motivates the verifier-grounded cross-check shown in Fig.~\ref{fig:fc_vs_judge_strength}: a Frama-C/WP implication test agrees with the judge on the subset it can decide, but leaves the majority of pairs incomparable under the per-goal budget.
\DIFaddend






\DIFaddbegin
\smallskip
\noindent\textbf{Frama-C/WP-based implication checking.}
We tested whether the LLM judge could be replaced by a direct Frama-C/WP implication check. For each both-valid pair we generated implication-checking harnesses for the postconditions and the matched loop invariants, and discharged the resulting obligations with Alt-Ergo and Z3 under a 5-second per-goal budget. On the \texttt{gpt-5.4-mini} both-valid set against \autospec, WP produced a definite ordering for only 51 of 242 pairs (21.1\%), leaving 191 pairs (78.9\%) incomparable. Among the WP-decided cases the direction agrees with the LLM judge, \toolname\ is stronger in 33 cases and \autospec\ in 10, in line with the judge's 47.9\% vs.\ 31.4\% ranking, which supports the judge on the subset WP can decide.
\DIFaddend

\DIFaddbegin
The large incomparable fraction should not be read as direct evidence of weak generated specifications. It mainly reflects the difficulty of automatically proving implication or equivalence between independently written specifications, a limitation also highlighted by the CLEVER benchmark~\cite{thakur2025clever}. For implication or equivalence checking of specifications, the verifier must discharge a conjunction of obligations: \texttt{requires}-implication, \texttt{ensures}-implication, and mutual loop-invariant obligations over typed memory, address shifts, recursive logic helpers, integer casts, and bit-vector operators, the same machinery that dominates the timeouts reported in Section~\ref{subsec:rq8-failure}. We found that raising the budget to several minutes moves only a handful of cases out of the incomparable category, while rendering the experiment intractable at this scale. We therefore use the LLM judge as the primary strength oracle in Table~\ref{tab:strength_summary}: it returns a verdict on every pair within a fixed budget, agrees in direction with WP wherever WP can decide, and remains usable across the JML/ACSL boundary where no single verifier applies.
\DIFaddend

\begin{figure*}[!t]
\centering
\begin{minipage}[t]{0.47\textwidth}
\centering
\resizebox{0.95\linewidth}{!}{%
\begin{tikzpicture}[x=0.019cm,y=0.45cm,font=\footnotesize,every node/.style={text=figgraytext}]
  \path[use as bounding box] (-100,-0.25) rectangle (275,4.35);
  \def\Wbar{242}
  \def\bh{0.24}
  \foreach \x in {0,50,100,150,200,242}{
    \draw[figgrid] (\x,0.45) -- (\x,4.10);
  }
  \foreach \y/\v/\p/\fillc/\linec/\lab in {%
    3.75/116/47.9/figcontractfill/figcontractstroke/{\toolname\ stronger},%
    2.80/76/31.4/figgrayfill/figgrayframe/{\autospec\ stronger},%
    1.85/26/10.7/figloopfill/figloopstroke/{Equal},%
    0.90/24/9.9/figgray!10/figgrayframe/{Incomparable}%
  }{
    \node[anchor=east] at (-3,\y) {\lab};
    \path[figbar,fill=\fillc] (0,\y-\bh) rectangle (\v,\y+\bh);
    \draw[figbarframe,draw=\linec] (0,\y-\bh) rectangle (\v,\y+\bh);
    \node[anchor=west] at ({\v+3},\y) {\v\ (\p\%)};
  }
  \draw[figaxis] (0,0.45) -- (\Wbar,0.45);
  \foreach \x in {0,50,100,150,200,242}{
    \draw[figaxis] (\x,0.45) -- (\x,0.30);
    \node[anchor=north] at (\x,0.25) {\x};
  }
\end{tikzpicture}
}%
\par\smallskip{\footnotesize\normalfont (a) LLM judge.}
\end{minipage}
\hfill
\begin{minipage}[t]{0.47\textwidth}
\centering
\resizebox{0.95\linewidth}{!}{%
\begin{tikzpicture}[x=0.019cm,y=0.45cm,font=\footnotesize,every node/.style={text=figgraytext}]
  \path[use as bounding box] (-100,-0.25) rectangle (275,4.35);
  \def\Wbar{242}
  \def\bh{0.24}
  \foreach \x in {0,50,100,150,200,242}{
    \draw[figgrid] (\x,0.45) -- (\x,4.10);
  }
  \foreach \y/\v/\p/\fillc/\linec/\lab in {%
    3.75/33/13.6/figcontractfill/figcontractstroke/{\toolname\ stronger},%
    2.80/10/4.1/figgrayfill/figgrayframe/{\autospec\ stronger},%
    1.85/8/3.3/figloopfill/figloopstroke/{Equal},%
    0.90/191/78.9/figgray!10/figgrayframe/{Incomparable}%
  }{
    \node[anchor=east] at (-3,\y) {\lab};
    \path[figbar,fill=\fillc] (0,\y-\bh) rectangle (\v,\y+\bh);
    \draw[figbarframe,draw=\linec] (0,\y-\bh) rectangle (\v,\y+\bh);
    \node[anchor=west] at ({\v+3},\y) {\v\ (\p\%)};
  }
  \draw[figaxis] (0,0.45) -- (\Wbar,0.45);
  \foreach \x in {0,50,100,150,200,242}{
    \draw[figaxis] (\x,0.45) -- (\x,0.30);
    \node[anchor=north] at (\x,0.25) {\x};
  }
\end{tikzpicture}
}%
\par\smallskip{\footnotesize\normalfont (b) Frama-C/WP implication check.}
\end{minipage}
\caption{Pairwise strength distribution on the \(242\) both-valid \texttt{gpt-5.4-mini} pairs against \autospec, comparing the LLM judge with a Frama-C/WP implication check.}
\label{fig:fc_vs_judge_strength}
\end{figure*}

\begin{finding}
\textbf{Finding (RQ6).} \toolname's specifications are judged stronger than \autospec's ($50.0\%$ vs.\ $31.5\%$) and on par with \specgen's; its strength advantage over \autospec\ grows with model capability.
\end{finding}
